\section{Basic Attention: de la compresión fija a la lectura según la necesidad}

Antes de llegar a self-attention, conviene dominar la attention de codificador-decodificador. El codificador produce un conjunto de estados $h_1,\ldots,h_n$, uno por cada posición de origen. En el instante $t$, el decoder usa el estado $s_{t-1}$ para calcular la compatibilidad, los pesos normalizados y el context:
\[
e_{t,i}=v^\top\tanh(W_hh_i+W_ss_{t-1}),\qquad
\alpha_{t,i}=\frac{\exp(e_{t,i})}{\sum_j\exp(e_{t,j})},\qquad
c_t=\sum_i\alpha_{t,i}h_i.
\]
Aquí $e_{t,i}$ es el score entre la $t$-ésima query y la posición de origen $i$, $\alpha_{t,i}$ es el attention weight normalizado y $c_t$ es el resultado de la lectura en este paso. Esta additive attention al estilo de Bahdanau sustituye el context vector único por un context distinto en cada paso.

\subsection{Soft attention y hard attention}

La \term{soft attention} asigna pesos continuos a todas las posiciones candidatas y, por lo general, puede entrenarse de extremo a extremo mediante retropropagación. La \term{hard attention} selecciona o muestrea unas cuantas posiciones; el cálculo puede ser más disperso, pero la selección discreta normalmente no es diferenciable de forma directa y requiere policy-gradient, relaxation u otro estimator. Entre ambas no hay un orden fijo del tipo «la soft es más precisa y la hard es más rápida»: el resultado real depende de la tarea y del método de entrenamiento.

\lecturefigure{slide-09-part-soft-attention.jpg}{La attention visual temprana ya distinguía entre soft selection y hard selection.}{Video oficial 06:30; cita de Xu et al. 2015 en la diapositiva.}

\begin{knowledgebox}{Lectura de la figura: los pesos continuos y la selección discreta cambian la interfaz de entrenamiento}
La soft attention equivale a un promedio ponderado diferenciable sobre varias regiones, de modo que el gradiente llega a cada posición; la hard attention equivale a elegir una región o unas pocas, con una ruta de cálculo más definida, pero introduce varianza de muestreo y el problema de asignación de crédito. El ejemplo del gato en la figura no busca demostrar que el modelo «entiende al gato», sino mostrar que los pesos de atención pueden funcionar como variable intermedia para encaminar la información. Además, visualizar los pesos no equivale de manera automática a una explicación causal.
\end{knowledgebox}

\begin{warningbox}{El attention weight no tiene por qué coincidir con una explicación humana}
Un peso alto indica en qué medida el cálculo actual lee cierto value, pero el modelo también pasa por residual, multi-head, nonlinearity y las capas posteriores. Un solo heatmap no basta para afirmar cuáles son las razones completas del razonamiento del modelo; hace falta verificarlo con ablation, counterfactual o causal intervention.
\end{warningbox}

\subsection{Global attention y local attention}

La \term{global attention} asigna un score a todas las posiciones de origen en cada decoding step; la \term{local attention} solo asigna scores dentro de una ventana predicha o de un vecindario seleccionado. La global ofrece acceso completo, pero su costo crece con la longitud; la local introduce un sesgo hacia los vecinos cercanos y reduce el cálculo, aunque puede pasar por alto evidencia lejana. Las propuestas posteriores de sliding-window, block-sparse y routing attention pueden verse como extensiones de este espacio de diseño.

\lecturefigure{slide-10-local-global-attention.jpg}{La global attention lee toda la secuencia; la local attention solo lee una ventana local.}{Video oficial 07:10; Luong et al. 2015.}

\begin{knowledgebox}{Lectura de la figura: la ventana es un sesgo inductivo, no solo un recurso de optimización}
El global model supone que cualquier posición puede ser relevante; el local model supone que la mayoría de las dependencias se concentran en el vecindario. El tamaño de la ventana, la forma de predecir su centro y el manejo de los bordes cambian las relaciones que el modelo puede expresar. En documentos largos, audio o video, la estructura local suele existir de verdad, pero las referencias entre segmentos son igual de importantes; el sistema debe elegir entre rendimiento, memoria y las dependencias de la tarea, en lugar de tratar la sparse attention como un sustituto sin pérdidas.
\end{knowledgebox}

\begin{importantbox}{La primera reestructuración de la interfaz con attention}
La attention temprana todavía conservaba el recurrent encoder/decoder, pero cambió la forma de leer: de un fixed vector a una content-dependent weighted retrieval. El siguiente paso, la self-attention, consiste en que cada posición de una misma secuencia lea las demás posiciones mediante una interfaz similar.
\end{importantbox}

\subsection{Resumen de la sección}

La basic attention ya contiene los ejes de diseño centrales de los sistemas posteriores: continua o discreta, global o local, acceso completo o encaminamiento disperso. La innovación del Transformer no consiste en inventar los «pesos», sino en convertir la attention en la capa principal de cálculo sobre secuencias.
